\documentclass[11pt,a4paper]{article}
\usepackage[margin=2.4cm]{geometry}
\usepackage{amsmath,amssymb,booktabs,hyperref}
\newcommand{\Lc}{\textbf{(c)}}
\newcommand{\Ld}{\textbf{(d)}}
\newcommand{\La}{\textbf{(a)}}
\newcommand{\scr}{\texttt{audit\_\allowbreak charges\_\allowbreak action.py}}
\emergencystretch=2em
\title{Independent audit of one-loop charges and static defect absorption}
\author{Project audit for Gustav Delius}
\date{4 November 2026}
\begin{document}
\maketitle
\begin{abstract}
The proposed collective-coordinate momentum is an exact derivative of the
proposed one-loop Lagrangian; a regulated Gaussian mode sum supports that
Lagrangian conditionally. \Lc\ Staticity does not justify changing which
bulk soliton charge enters absorption. The defect's nonlinear boundary
potential and localized charge operators require a renormalization argument
beyond the bulk tadpole, which vanishes on its constant vacua. \Ld\ A shift
of the absorption locus reconciles both charge pairs algebraically \Lc,
but no quantum sewing calculation establishes it physically \Ld.
Existing good-equation conventions and bulk subtraction are retained;
no new renormalization prescription is chosen.
\end{abstract}

\section{Scope, labels, and evidence}
Labels are \La\ literature results, \textbf{(b)} literature heuristics,
\Lc\ project derivations, and \Ld\ unestablished physical identifications or
possibilities. Every displayed new \Lc\ identity is checked independently in
\texttt{phase2/scripts/audit\_charges\_action.py}, using SymPy and 65-digit
mpmath. The script imports no existing project charge formulas.

The main audited passages are \texttt{charge\_frame\_master.tex}, lines
284--316 (co-moving action), 325--363 (derivative), and 367--398 (static-defect
interpretation); also \texttt{oneloop\_\allowbreak defect\_\allowbreak phase.tex}, lines 605--650
(renormalization and absorption). The previous independent Phase 2h audit was
consulted as a lead; the checks here were implemented separately.

\section{What the co-moving mode sum does establish}
Write the potential fluctuation in $\xi=x-vt$ as $\varphi$, and let
$u_0=\partial_\xi\phi_0$. Taylor expansion of the fixed project action gives
\Lc\ (\scr)
\begin{equation}
 \mathcal L_2=\tfrac12(\varphi_t-v\varphi_\xi)^2
 -\tfrac12(1-6u_0)\varphi_\xi^2-\tfrac12\varphi_{\xi\xi}^2.
 \label{eq:quadratic}
\end{equation}
Thus asymptotic continuum modes have co-moving frequency
$\Omega(q)=|q|\sqrt{1+q^2}-vq$. This follows by substitution of a plane wave
into the vacuum quadratic action, and is positive for $q\ne0$, $|v|<1$.
The determinant interpretation still requires a complete, correctly normalized
mode basis, the collective-coordinate projection, and the specified boundary
conditions. This audit does not supply those parts. \Ld

There is nevertheless a useful conditional derivation beyond checking a
Legendre identity. If equal-mode-count Gaussian quantization is valid, its
continuum contribution uses the derivative form of the spectral trace,
including its cutoff endpoint. For $\theta(q)=\sum_r c_r\Theta(q;\phi_r)$
on the existing soliton branches, \Lc\ (\scr):
\begin{align}
 E_{\rm dens}(\Lambda)&=\frac1{4\pi}\int_{-\Lambda}^{\Lambda}
          |q|\sqrt{1+q^2}\,\theta'(q)\,dq,\\
 B(\Lambda)&=\frac{\Lambda\sqrt{1+\Lambda^2}}{4\pi}
                   [\theta(\Lambda)-\theta(-\Lambda)],\\
 E_{\rm mode}(\Lambda)&=E_{\rm dens}(\Lambda)-B(\Lambda),\\
 \Delta E_{\rm ren}&=\lim_{\Lambda\to\infty}
 \left[E_{\rm mode}(\Lambda)+\frac{3k}{\pi}
                   (\sqrt{1+\Lambda^2}-1)\right].
 \label{eq:regulated}
\end{align}
The final line reproduces the existing closed energy formula; it does not
change its finite part. The endpoint satisfies
$B(\Lambda)/\Lambda\to3k/(2\pi)$, so it cannot be omitted before applying
the existing counterterm. The momentum endpoint
$[q\theta(q)]_{-\Lambda}^{\Lambda}/(4\pi)$ tends to zero. Consequently
the same conditional prescription gives
\begin{equation}
 L_1(v)=v\Delta P_{\rm mode}-\Delta E_{\rm ren}.
 \label{eq:conditionalL}
\end{equation}
This is \Lc\ as a consequence of the stated spectral assumptions, rather
than a completed physical determinant derivation. The original
\texttt{charge\_and\_frame\_audit.py:96--99} defines this combination;
its quadrature at lines 220--245 checks that definition and the already
subtracted energy integrand.

The equal-mode-count endpoint is also explicitly discussed by Faddeev and
Korepin in Appendix 6, equations A.6.11--A.6.14. \La\ Only the relevant
full-text passages, rather than the whole paper, were read for this audit.

\section{When differentiation gives the physical momentum}
For an on-shell effective action with valid boundary variation,
differentiation with respect to velocity computes the canonical charge. The
first-order correction includes the classical charge evaluated on the
corrected profile. It need not equal a fixed-background expectation alone.
The finite-dimensional version makes this precise: let
$L=L_0(v,y)+\hbar L_1(v,y)$, $L_{0,y}(v,y_0)=0$, and
$y=y_0+\hbar y_1$. Assuming a nondegenerate Hessian after zero-mode projection,
\Lc\ (\scr)
\begin{equation}
 y_0'=-\frac{L_{0,yv}}{L_{0,yy}},\qquad
 y_1=-\frac{L_{1,y}}{L_{0,yy}},\qquad
 P_1=L_{1,v}+L_{0,vy}y_1
     =\frac{d}{dv}L_1(v,y_0(v)).
 \label{eq:saddle}
\end{equation}
The field-theory application requires the corresponding functional inverse
and boundary terms. In particular, changing the Boussinesq velocity also
changes the asymptotic potential jump. The finite-dimensional identity does
not check those boundary conditions. The physical identification of the
published $\Delta P_{\rm can}$ is therefore conditional, although its
derivative formula is correct. \Ld

Taking \eqref{eq:conditionalL} as input yields, independently at 65 digits,
\Lc\ (\scr)
\begin{align}
 \Delta P_{\rm can}&=\frac{dL_1}{dv}
 =-\frac{\sqrt3\cos\alpha}{\pi}
   \log\left|\frac{\tan((\alpha+2\pi/3)/2)}
                    {\tan((\alpha-2\pi/3)/2)}\right|,\quad (\epsilon=+1),\\
 \Delta E_{\rm can}&=v\Delta P_{\rm can}-L_1,\qquad
 \frac{d\Delta E_{\rm can}}{d\alpha}
       =v\frac{d\Delta P_{\rm can}}{d\alpha}.
\end{align}
These equations prove internal consistency, not a charge Ward identity for
the renormalized field theory. Also, the printed intermediate equation
for $-dL_1/d\alpha$ in the master note has the wrong overall sign:
\Lc\ (\scr)
\begin{equation}
 \frac{dL_1}{d\alpha}=\frac1{2\pi}\sum_r c_r
 [\cos(2\phi_r)-\cos(\alpha+\phi_r)]
 \log|\tan(\phi_r/2)|.
\end{equation}
The light-root cancellation survives this correction.

\section{A static defect still has a localized momentum operator}
Corrigan--Zambon give the conserved modified momentum, including localized
field terms, in equations 3.11 and 3.49--3.51. It is not the ordinary
translation Noether momentum of a translation-invariant system: a fixed
defect breaks that symmetry. Their $H$ contains a fifth-degree potential
in the field jump (3.42--3.43), and their $h$ is fourth-degree (3.46).
\La\ Relevant full-text sections 2--3 and 5.1 were read online, not the
whole article. The dictionary is $U_{\rm CZ}=-3\phi_L$,
$V_{\rm CZ}=-3\phi_R$, and $\mathcal L_{\rm CZ}=9\mathcal L$;
$c=\lambda=\kappa=1$.

For direct independent verification, write $A=\eta/\sqrt3$, $B=2A$,
and denote a constant vacuum by $\sigma$. Expanding the boundary charges
about that vacuum in project-normalized potential fluctuations
$a,b$ gives \Lc\ (\scr, displayed for $\eta=+1$, with the same algebra
for the other sign)
\begin{align}
 P_{\mathcal D,2}
 &=-\frac A2(a_x^2+b_x^2)-B a_xb_x
       -\frac A4(1-\sigma^2)(a-b)^2,\label{eq:P2}\\
 H_2&=\sigma\left[\frac B2(a_x^2+b_x^2)+A a_xb_x\right]
   +\frac{1-\sigma^2}{2}(a-b)(a_x+b_x)\nonumber\\
   &\quad-\frac{B\sigma(1-\sigma^2)}8(a-b)^2,
   \qquad E_{\mathcal D,2}=-H_2.
\end{align}
They are generally nonzero although the defect has no translational
collective coordinate. A complete spectral trace of the conserved charge
might already include these terms through its mode normalization. This audit
does not assert it cannot. What is absent from the published check is a
derivation equating that trace to $\int q\theta_D'(q)\,dq/(4\pi)$.
Staticity alone cannot supply that derivation. \Ld

\section{The defect subtraction requires a boundary derivation}
For the existing two-root kernel, put $\cos\gamma=\sigma/2$,
$\psi_\pm=\pm2\pi/3-\gamma$ and weights $(-\eta,+\eta)$.
Then \Lc\ (\scr)
\begin{equation}
 S_D:=\sum c_D\sin\psi=-\eta\sqrt3\cos\gamma
          =-\frac{\eta\sqrt3\sigma}{2},\qquad
 -\frac{\omega'(q)[\theta_D(q)-\theta_D(-q)]}{4\pi}
                 \longrightarrow\frac{S_D}{\pi}\quad(q\to+\infty).
\end{equation}
For example $\eta=+1$, $\sigma=-1.5$ gives slope
$0.413496671566\ldots$. This is a nonzero continuum divergence even on
a linearly stable vacuum.

The existing bulk normal-ordering counterterm is proportional to the sum of
$\int\phi_x\,dx$ over the two open half-lines. Each integrand vanishes
on a constant defect vacuum. Hence that bulk term cannot by itself cancel
the displayed divergence. Interpreting the jump as a distribution at the
defect would introduce a boundary prescription which must be derived and
specified. The statement in \texttt{oneloop\_defect\_phase.tex:611--613}
that the potential is linear is contradicted by its fifth-degree term.
Thus the finite defect energy formula is an algebraically coherent subtraction,
but its physical renormalization argument remains unestablished. \Ld
No replacement prescription is adopted here.

\section{A conditional absorption-locus shift}
Suppose both that the canonical soliton charges are physical and that the
quoted defect vacuum shifts are physical, while retaining their algebraic
mode-charge jump identity. Conservation would then require more than the
published assertion that static absorption uses a different bulk momentum.
There is a simple kinematic possibility: the absorption soliton could have
$\alpha_{\rm abs}=\alpha_0+\hbar a_1+O(\hbar^2)$.
Because \Lc\ (\scr)
\begin{equation}
 \Delta E_{\rm can}-\Delta E_{\rm ren}
       =v(\Delta P_{\rm can}-\Delta P_{\rm mode}),\qquad
 E_{\rm cl}'=vP_{\rm cl}',
\end{equation}
the choice
\begin{equation}
 a_1=\frac{\Delta P_{\rm mode}-\Delta P_{\rm can}}{P_{\rm cl}'}
 \label{eq:shift}
\end{equation}
makes both first-order canonical soliton charges equal the mode-charge
increments, whenever $P_{\rm cl}'\ne0$. This is a conditional algebraic
observation \Lc, not a derived quantum absorption condition \Ld.
At $\alpha_0=0.15,0.5,0.85$ the coefficients are respectively
$0.753417205691$, $0.920879681850$, and $2.33030395807$.
Quantum shifts of defect parameters or additional localized charge terms
are also not excluded. Selecting among these requires a renormalized
quantum sewing or charge calculation, not a choice made in this audit.

\section{Checks, limitations, and failures recorded}
\begin{center}\small
\begin{tabular}{p{.53\linewidth}p{.36\linewidth}}
\toprule
Claim checked & Verification\\\midrule
Co-moving quadratic action & Exact SymPy Taylor coefficient\\
Existing bulk energy with cutoff endpoint & 65-digit closed integration;
cutoff $10^{20}$, residual below $10^{-19}$\\
Legendre formulas and absorption-shift identities & Derivatives checked to
$10^{-59}$ or better at three regular light solitons\\
Stationary profile correction & Exact symbolic chain-rule identity\\
Localized quadratic defect charges & Independent symbolic expansion\\
Defect UV slope & Exact trigonometry and 65-digit large-$q$ check\\\bottomrule
\end{tabular}
\end{center}
Failures recorded, including statements checked and found false:
\begin{enumerate}
\item The master note's profile at line 250 is
$-k^2\operatorname{sech}^2(k\xi/2)$, whereas the fixed convention gives
$+k^2\operatorname{sech}^2(k\xi/2)/2$. The displayed profile fails the
travelling-wave equation already at its centre (checked symbolically).
\item The weighted cosine sum stated to vanish at line 257 is actually
$\sum c_r\cos\phi_r=-3\cos\alpha$. The final charge derivative can still be
correct because cancellations occur within the energy derivative.
\item The sign of the printed $-dL_1/d\alpha$ weight formula fails;
the same weight equals $+dL_1/d\alpha$. The closed canonical momentum passes.
\item The defect potential is not linear. Its fifth-degree term and its
localized quadratic momentum contradict the reasons given for discarding
boundary effects.
\item No implementation assertion failed during development of this audit
script; no tolerance was weakened. A direct DOI webpage lookup for
Faddeev--Korepin was inaccessible; the DOI was subsequently verified against
the publisher's indexed entry.
\end{enumerate}
The full determinant, projected measure, quantum boundary counterterms,
and quantum charge Ward identities were not computed. The successful script
therefore ends with \texttt{ALL CHECKS PASSED} only for its stated algebraic
and conditional checks, not for the unresolved physical identifications.

\begin{thebibliography}{9}
\bibitem{CZ} E.~Corrigan and C.~Zambon, \emph{Adding integrable defects to the
Boussinesq equation}, \href{https://arxiv.org/abs/2306.03772}{arXiv:2306.03772}.
Relevant full-text sections 2--3 and 5.1 read online for this audit.
\bibitem{FK} L.~D.~Faddeev and V.~E.~Korepin, \emph{Quantum theory of solitons},
Physics Reports \textbf{42} (1978), 1--87,
\href{https://doi.org/10.1016/0370-1573(78)90058-3}{doi:10.1016/0370-1573(78)90058-3}.
DOI verified against the publisher's indexed entry; reading scope stated above.
\end{thebibliography}
\end{document}
